\documentclass[11pt]{article}

\usepackage[english]{babel}
\usepackage{kotex}    

\usepackage[letterpaper,top=1.2cm,bottom=1.2cm,left=1.2cm,right=1.2cm,marginparwidth=1.75cm]{geometry}

% Useful packages
\usepackage{amsmath, amsthm, amssymb, graphicx, enumitem, authblk, tikz, tikz-cd, verbatim, relsize}
\usepackage{thmtools}
\usepackage{tcolorbox}
\usepackage{forest} % 트리 구조를 위한 패키지
\usepackage{mathtools}
\usepackage[colorlinks=true, allcolors=black]{hyperref}
\usepackage[framemethod=TikZ]{mdframed}
\usepackage{lmodern}

\DeclareMathAlphabet{\mathtt}{OT1}{lmtt}{m}{n}

\usetikzlibrary{calc}
\usetikzlibrary{shapes,backgrounds}
\usetikzlibrary{decorations.pathreplacing}
\usetikzlibrary{arrows.meta, positioning, matrix, fit}

% \usetikzlibrary{arrows.meta, positioning, external}
% \pgfplotsset{compat=1.18} % Set compatibility for pgfplots

\newtheoremstyle{romanstyle} % 스타일 이름
  {10pt} % 위쪽 여백
  {10pt} % 아래쪽 여백
  {\rmfamily} % 본문 로마체
  {} % 들여쓰기 없음
  {\bfseries} % 제목 굵게
  {.} % 제목 뒤 구두점
  { } % 제목과 본문 사이의 기본 간격
  {} % 제목 레이아웃 기본값

% 새 스타일 적용
\theoremstyle{romanstyle}


% 1) 박스형 스타일 프리셋
\declaretheoremstyle[
  headfont=\bfseries, bodyfont=\rmfamily,
  headpunct={.}, postheadspace=.5em,
  mdframed={linewidth=0.6pt, linecolor=black, backgroundcolor=white,
            innerleftmargin=8pt, innerrightmargin=8pt,
            innertopmargin=6pt, innerbottommargin=6pt,
            skipabove=10pt, skipbelow=10pt, roundcorner=0pt}
]{thm-plain}

\declaretheoremstyle[
  headfont=\bfseries, bodyfont=\rmfamily, headpunct={.},
  postheadspace=.5em,
  mdframed={linewidth=1pt, linecolor=black, backgroundcolor=gray!4,
            leftline=true, rightline=false, topline=false, bottomline=false,
            innerleftmargin=10pt, innerrightmargin=8pt,
            innertopmargin=6pt, innerbottommargin=6pt,
            skipabove=10pt, skipbelow=10pt}
]{thm-left}

\declaretheoremstyle[
  headfont=\bfseries, bodyfont=\rmfamily, headpunct={.},
  postheadspace=.5em,
  mdframed={linewidth=0.8pt, linecolor=black, backgroundcolor=white,
            roundcorner=2pt, tikzsetting={dash pattern=on 3pt off 2pt},
            innerleftmargin=8pt, innerrightmargin=8pt,
            innertopmargin=6pt, innerbottommargin=6pt}
]{thm-dashed}

\declaretheoremstyle[
  headfont=\bfseries, bodyfont=\rmfamily, headpunct={.},
  postheadspace=.5em,
  mdframed={linewidth=0.8pt, linecolor=blue!60!black, backgroundcolor=black!3,
            roundcorner=1.5pt,
            innerleftmargin=8pt, innerrightmargin=8pt,
            innertopmargin=6pt, innerbottommargin=6pt}
]{thm-blue}

\declaretheorem[style=thm-left,   numberwithin=subsection]{theorem}
\declaretheorem[style=thm-blue,   numberwithin=subsection]{proposition}
\declaretheorem[style=thm-dashed, numberwithin=subsection]{lemma}
\declaretheorem[style=thm-blue,   numberwithin=subsection]{corollary}
\declaretheorem[style=thm-plain,  numberwithin=subsection]{definition}

\newtheorem{axiom}{Axiom}

% Unnumbered environments
\newtheorem*{observation}{Observation}
\newtheorem*{example}{Example}
\newtheorem*{solution}{Solution}

% Mathematical commands
\newcommand{\st}{\text{ s.t. }}
\newcommand{\as}{\text{ as }}
\newcommand{\setx}{\{x_0, x_1, \dots, x_n\}}
\newcommand{\setn}{\{1, \dots, n\}}
\newcommand{\dis}{\displaystyle}
\newcommand{\lete}{Let $\epsilon > 0$ be given.}
\newcommand{\fs}{\text{ for some }}
\newcommand{\defequal}{\overset{\textnormal{def}}{=}}
\newcommand{\setequal}{\overset{\textnormal{set}}{=}}
\newcommand*{\aand}{_{\text{ and }}}
\newcommand*{\oor}{_{\text{ or }}}
\DeclareMathOperator{\argmax}{argmax}
\DeclareMathOperator{\argmin}{argmin}
\newcommand{\im}{\operatorname{im}}
\newcommand{\nullity}{\operatorname{nullity}}
\newcommand{\rank}{\operatorname{rank}}
\newcommand{\ipvs}{\langle \, \cdot \, \rangle}
\newcommand{\diam}{\operatorname{diam}} % Add this line in the preamble
\newlength{\datedwd}
\newcommand{\dated}[2]{%
  \settowidth{\datedwd}{#1}%
  \noindent
  \parbox[t]{\datedwd}{\strut #1}%
  \hspace{0.5em}%
  \begin{minipage}[t]{\dimexpr\linewidth-\datedwd-0.5em\relax}
    \vspace{-0.65em}%
    \begin{enumerate}[label=\arabic*., leftmargin=*, itemsep=0pt, topsep=0pt, parsep=0pt]
      #2
    \end{enumerate}
  \end{minipage}\par
}

% ===================== TOC에 정리 제목(옵션 있을 때만) 자동 추가 ===================== %
\usepackage{xparse}   %% <<< ADDED
\usepackage{etoolbox} %% <<< ADDED
\setcounter{tocdepth}{3}    %% <<< ADDED: subsubsection까지 TOC 표시
\setcounter{secnumdepth}{3} %% <<< ADDED: (선택) 번호도 3단계까지

% theorem
\let\theoremOLD\theorem           %% <<< ADDED
\let\endtheoremOLD\endtheorem     %% <<< ADDED
\RenewDocumentEnvironment{theorem}{o}{%% <<< ADDED
  \IfNoValueTF{#1}{%
    \theoremOLD
  }{%
    \theoremOLD
    \phantomsection
    \addcontentsline{toc}{subsubsection}{%
      \protect\numberline{\thetheorem}\ #1}%
  }%
}{\endtheoremOLD}

% lemma
\let\lemmaOLD\lemma
\let\endlemmaOLD\endlemma
\RenewDocumentEnvironment{lemma}{o}{%
  \IfNoValueTF{#1}{%
    \lemmaOLD
  }{%
    \lemmaOLD
    \phantomsection
    \addcontentsline{toc}{subsubsection}{%
      \protect\numberline{\thelemma}\ #1}%
  }%
}{\endlemmaOLD}

% proposition
\let\propositionOLD\proposition
\let\endpropositionOLD\endproposition
\RenewDocumentEnvironment{proposition}{o}{%
  \IfNoValueTF{#1}{%
    \propositionOLD
  }{%
    \propositionOLD
    \phantomsection
    \addcontentsline{toc}{subsubsection}{%
      \protect\numberline{\theproposition}\ #1}%
  }%
}{\endpropositionOLD}

% corollary
\let\corollaryOLD\corollary
\let\endcorollaryOLD\endcorollary
\RenewDocumentEnvironment{corollary}{o}{%
  \IfNoValueTF{#1}{%
    \corollaryOLD
  }{%
    \corollaryOLD
    \phantomsection
    \addcontentsline{toc}{subsubsection}{%
      \protect\numberline{\thecorollary}\ #1}%
  }%
}{\endcorollaryOLD}

% definition
\let\definitionOLD\definition
\let\enddefinitionOLD\enddefinition
\RenewDocumentEnvironment{definition}{o}{%
  \IfNoValueTF{#1}{%
    \definitionOLD
  }{%
    \definitionOLD
    \phantomsection
    \addcontentsline{toc}{subsubsection}{%
      \protect\numberline{\thedefinition}Definition\ #1}%
  }%
}{\enddefinitionOLD}
% ==================================================================== %

\title{연구인턴장학 결과보고서}
\author{정종원}
\affil{서울시립대학교 수학과}
\date{}
% \sloppy
\begin{document}
\maketitle


\section{Preliminaries}

\subsection{Hilbert Space}
\begin{definition}
    Complete Inner product Vector Space is called \textit{Hilbert Space}.
\end{definition}

\subsection{Hilbert Space in $\mathbb{R}^\omega$}
\begin{definition}
    Define $\dis \mathbb{R}^\omega \defequal \prod_{i=1}^\infty \mathbb{R}$ as the countable product of Euclidean space $\mathbb{R}$ with product topology. \\
    And define $\dis \mathbb{H} \defequal \left\{ {\{x_n \}}_{n=1}^\infty \;\middle|\; \sum_{n=1}^{\infty} x^2_n < \infty \right\} \subset \mathbb{R}^\omega$, 
    \textit{Metric} on $\mathbb{H}$ as $\dis \mu:\mathbb{H} \times \mathbb{H} \to \mathbb{R}: ( \{x_n\}, \{y_n\}) \mapsto \sqrt{\sum_{i=1}^{\infty} (x_i - y_i)^2}$.\\
    The Metric Space $(\mathbb{H}, \mu)$ is called \textit{Hilbert Space} or \textit{$l_2$ Space}.    
\end{definition}
Define the operations elementwise; then $(\mathbb{H}, +, \times)$ is a Vector Space over $\mathbb{R}$.\\
Moreover, $\mathbb{H}$ is Complete Metric Space and Inner product Vector Space.
\begin{lemma}
    $\dis \mu:\mathbb{H} \times \mathbb{H} \to \mathbb{R}: ( \{x_n\}, \{y_n\}) \mapsto \sqrt{\sum_{i=1}^{\infty} (x_i - y_i)^2}$
    is Metric function induced by the inner product.
\end{lemma}
    \begin{proof}
    We know that $\mathbb{R}^\omega$ is Vector Space. Moreover, $\mathbb{H} \subset \mathbb{R}^\omega$ is Subspace.
    Using subspace criteria:
    \begin{center}
        $S \subset V$ is Subspace of Vector Space $V$ if and only if $0 \in S$ and For any $x,y \in S$ and $a \in F$, $ax + y \in S$.
    \end{center}
    Clearly, $\{0\} \in \mathbb{H}$. Let $a \in \mathbb{R}$ and $\{x_n\}, \{y_n\} \in \mathbb{H}$ be given. Then, $a \{x_n\} + \{y_n\} = \{ax_n + y_n\} \in \mathbb{H}$ because:
    \begin{align*}
        \sum_{i=1}^{\infty} (a x_i + y_i)^2
        = \sum_{i=1}^{\infty} \left[a^2 x_i^2 + 2a x_i y_i + y_i^2\right]
        \overset{(*)}{=} a^2 \sum_{i=1}^{\infty} x_i^2 + 2a \sum_{i=1}^{\infty} x_i y_i + \sum_{i=1}^{\infty} y_i^2 < \infty
    \end{align*}
    The $(*)$ given by:
    \begin{align*}
        \sum_{i=1}^{\infty} |x_i y_i|
        =\sum_{i=1}^{\infty} |x_i| |y_i|
        \leq \sum_{i=1}^{\infty} (\max(|x_i|, |y_i|))^2
        \leq \sum_{i=1}^{\infty} (x_n^2 + y_n^2)
        = \sum_{i=1}^{\infty} x_n^2 + \sum_{i=1}^{\infty} y_n^2 < \infty
        \tag{$*$}
    \end{align*}
    Thus $\mathbb{H}$ is Vector Space over $\mathbb{R}$. Now, define \textit{inner product} on $\mathbb{H}$ as:
    \begin{align*}
        \langle \ \cdot \ , \ \cdot \ \rangle : \mathbb{H} \times \mathbb{H} \to \mathbb{R}:
        \left(\{x_n\}, \{y_n\} \right) \mapsto \sum_{i=1}^{\infty} x_i y_i
    \end{align*}
    This definition is well-defined since $(*)$. And, Linearity in first:
    \begin{align*}
        \langle a \{x_n\} + \{y_n\} , \{z_n\} \rangle
        = \langle \{a x_n + y_n\} , \{z_n\} \rangle
        = \sum_{i=1}^{\infty} (ax_i + y_i) z_i
        = a \sum_{i=1}^{\infty} x_i z_i + \sum_{i=1}^{\infty} y_i z_i
        = a \langle \{x_n\}, \{z_n\} \rangle + \langle \{y_n\} , \{z_n\} \rangle
    \end{align*}
    The other conditions are clear. Thus, $(\mathbb{H}, \langle \cdot , \cdot \rangle)$ is \textit{inner product space}. \\
    Using \textit{ineer product}, define the \textit{Norm} on $\mathbb{H}$ as:
    \begin{align*}
        \| \cdot \| : \mathbb{H} \to \mathbb{R} : \{x_n\} \mapsto \sqrt{\langle \{x_n\}, \{x_n\} \rangle}
    \end{align*}
    Finally, define \textit{Metric} on $\mathbb{H}$ as:
    \begin{align*}
        \mu:\mathbb{H} \times \mathbb{H} \to \mathbb{R}: ( \{x_n\}, \{y_n\}) \mapsto \| \{x_n\} - \{y_n\} \| = \sqrt{\sum_{i=1}^{\infty} (x_i - y_i)^2}
    \end{align*}
    \end{proof}
\begin{theorem}
    \textit{Hilbert Space} is Separable.
\end{theorem}
\begin{proof}
    For each $n \in \mathbb{N}$, define $D_n \defequal \{\{p_n\} \mid p_i \in \mathbb{Q},\ p_{n+1} = p_{n+1} = \cdots = 0 \}$ and $\dis D \defequal \bigcup_{n \in \mathbb{N}} D_n$. \\
    Then, $D$ is countable set. We will show that $\overline{D} = \mathbb{H}$. \\
    Let $\epsilon > 0$ and $\{x_n\} \in \mathbb{H}$ be given. Since convergence, there exists $N \in \mathbb{N}$ such that 
    \begin{align*}
        \sum_{i=N+1}^{\infty} x^2_i = \sum_{i=1}^{\infty} x^2_i - \sum_{i=1}^{N} x^2_i < \frac{\epsilon^2}{2}
    \end{align*}
    Since density of Rationals, put each $i = 1, 2 , \ldots, N$, $p_i \in \mathbb{Q}$ $\dis |x_i - p_i| < \frac{\epsilon}{\sqrt{2N}}$ and $p_i = 0$ for $i \geq N+1$. \\
    Then, $\{p_n\} \in D_n \subset D$ and
    \begin{align*}
        \mu\left( \{x_n\}, \{p_n\} \right)
        = \sqrt{\sum_{i=1}^{N} (x_i - p_i)^2 + \sum_{i=N+1}^{\infty} (x_i - p_i)^2}
        = \sqrt{\sum_{i=1}^{N} (x_i - p_i)^2 + \sum_{i=N+1}^{\infty} x_i^2}
        < \sqrt{N \cdot \frac{\epsilon^2}{2N} + \frac{\epsilon^2}{2}}
        = \epsilon 
    \end{align*}
\end{proof}
\begin{corollary}
    \textit{Hilbert Space} is Second-Countable.
\end{corollary}

\begin{theorem}
    \textit{Hilbert Space} is Complete.
\end{theorem}
\begin{proof}
    Let $\{ \{x_{n,i} \}_{i=1}^\infty \}_{n=1}^\infty$ be a Cauchy sequence in $\mathbb{H}$. For any fixed $n,m \in \mathbb{N}$ and for each $j \in \mathbb{N}$,
    \begin{align*}
        |x_{n,j} - x_{m,j}| < \mu(\{x_{n,i}\}, \{x_{m,i}\})
        = \sqrt{\sum_{i=1}^{\infty} (x_{n,i} - x_{m,i})^2}
    \end{align*}
    That is, for each $j \in \mathbb{N}$, $\{x_{n,j}\}$ is Cauchy sequence in $\mathbb{R}$. Since $\mathbb{R}$ is Complete,
    put $\dis y_j \defequal \lim_{n \to \infty} x_{n,j}$, each $j \in \mathbb{N}$. \\
    Let $\epsilon > 0$ be given. Then, there exists $N \in \mathbb{N}$ such that $n,m \geq N \implies \dis \mu(\{x_{n,i}\}, \{x_{m,i}\}) < \frac{\epsilon}{2}$. \\
    Meanwhile, for each $k \in \mathbb{N}$,
    \begin{align*}
        \sum_{i=1}^{k} (x_{n,i} - x_{m,i})^2 \leq \sum_{i=1}^{\infty} (x_{n,i} - x_{m,i})^2 = \left[\mu(\{x_{n,i}\}, \{x_{m,i}\})\right]^2
    \end{align*}
    Thus, $n,m \geq N \implies \dis \sum_{i=1}^{k} (x_{n,i} - x_{m,i})^2 < \left( \frac{\epsilon}{2} \right)^2$, for each $ k \in \mathbb{N}$.\\
    Taking limit to $m$, then $n \geq N \implies \dis \lim_{m \to \infty} \left(\sum_{i=1}^{k} (x_{n,i} - x_{m,i})^2 \right) = \sum_{i=1}^{k} \left(x_{n,i} - \lim_{m \to \infty} x_{m,i}\right)^2 = \sum_{i=1}^{k} (x_{n,i} - y_i)^2 < \left(\frac{\epsilon}{2}\right)^2$.\\
    And, for all $k \in \mathbb{N}$,
    \begin{align*}
        \sum_{i=1}^{k} y_i^2 = \sum_{i=1}^{k} (2(x_{n,i}^2 + (x_{n,i} - y_i)^2)) \leq 2\|\{x_{n,i}\}_{i=1}^\infty \|^2 + \left(\frac{\epsilon}{2}\right)^2
    \end{align*}
    Thus $\{y_i\} \in \mathbb{H}$. As a result,
    \begin{align*}
        n \geq N \implies \mu(\{x_n\}, \{y_n\})
        = \sqrt{\sum_{i=1}^{\infty} (x_{n,i} - y_i)^2}
        = \sqrt{\lim_{k \to \infty} \sum_{i=1}^{k} (x_{n,i} - y_i)^2}
        < \frac{\epsilon}{2}
    \end{align*}
\end{proof}


\subsection{Vector Space}

\begin{definition}
    Suppose that $F$ is a Field, and $V$ is an Abelian Group. \\
    And, the operation $\cdot: F \times V \to V$ satisfies: For any $a, b \in F$ and $v, w \in V$,
    $\begin{cases}
        a \cdot (v + w) = a\cdot v + a\cdot w \\
        (a + b) \cdot v = a\cdot v + b\cdot v \\
        (ab) \cdot v = a \cdot (b \cdot v) \\
        1 \cdot v = v
    \end{cases}$ \\
    The triple $(V, +, \cdot)$ is called the \textit{Vector Space} over $F$.
\end{definition}
Equivalently, The Vector Space over a field $F$ is $F$-Module.
\begin{definition}
    Suppose that $V, W$ are Vector Space over a field $F$.\\
    A map $\mathcal{L}:V \to W$ is called \textit{Linear Map} if:
    \begin{center}
        For any $a \in F$ and $v_1, v_2 \in V$, $\mathcal{L}(a \cdot v_1 + v_2) = a \cdot \mathcal{L}(v_1) + \mathcal{L}(v_2)$
    \end{center}
\end{definition}
\begin{definition}
    Suppose that $V$ is a Vector Space over a field $F$, and $W \subseteq V$ is a Subset. \\
    The $W$ is called \textit{Subspace} of $V$ if:
    $\begin{cases}
        \text{For any $a \in F,\ w \in W,\ a \cdot w \in W$} \\
        \text{For any $w_1, w_2 \in W,\ w_1 + w_2 \in W$}
    \end{cases}$. \\
\end{definition}
That is, the Subset of a Vector Space which is closed under the addition and scalar multiplication, then it is a Vector Space.
\begin{lemma}
    Arbitrary intersection of Subspace is a Subspace.
\end{lemma}
\begin{proof}
    Suppose that $V$ is a Vector Space, and $W_\alpha \leq V,\ \alpha \in \Lambda$ are Subspaces. \\
    Using Subspace Criterion: Let $a \in F,\ w_1, w_2 \in \bigcap_{\alpha \in \Lambda} W_\alpha$ be given.
    Then, for any $\alpha \in \Lambda$, $a \cdot w_1 + w_2 \in W_\alpha$.
\end{proof}
This Lemma allows the definition:
\begin{definition}
    Suppose that $V$ is a Vector Space over a field $F$, and $S \subseteq V$ be a Subset. \\
    Define a \textit{Generated Subspace} by $S$ is:
    \begin{align*}
        \langle S \rangle \defequal \bigcap_{S \subseteq W \leq V} \hspace{-0.2cm} W
    \end{align*}
\end{definition}
This $\langle S \rangle$ is the \textit{unique smallest} Subspace containing $S$. This $S$ is called \textit{Generating Subset} of $\langle S \rangle$.
\begin{lemma}
    Suppose that $V$ is a Vector Space over a field $F$, and $S \subseteq V$ be a Subset. Then,
    \begin{align*}
        \langle S \rangle = \{ a_1 \cdot v_1 + \cdots + a_n \cdot v_n \mid n \in \mathbb{N},\ a_i \in F,\ v_i \in S \}
    \end{align*}
\end{lemma}
\begin{proof}
    First, $\langle S \rangle \supseteq \{ a_1 \cdot v_1 + \cdots + a_n \cdot v_n \mid n \in \mathbb{N},\ a_i \in F,\ v_i \in S \}$, because $\langle S \rangle$ is closed under the opeartions. \\
    And, the set $\{ a_1 \cdot v_1 + \cdots + a_n \cdot v_n \mid n \in \mathbb{N},\ a_i \in F,\ v_i \in S \}$ is a Subgroup of $V$ containing $S$, \\
    Hence $\langle S \rangle \subseteq \{ a_1 \cdot v_1 + \cdots + a_n \cdot v_n \mid n \in \mathbb{N},\ a_i \in F,\ v_i \in S \}$.
\end{proof}


\newpage

\section{Three-Body Problem}
\subsection{Newtonian Formulation}

\begin{definition}[Newtonian Three-Body Problem]
    Fix $G>0$. Let $m_1,m_2,m_3>0$ be masses and let
    \[
        \mathbf{r}_i:\mathbb{R}\to\mathbb{R}^3\qquad (i=1,2,3)
    \]
    be position functions. The \textit{(Newtonian) three-body problem} asks for trajectories
    $\{\mathbf{r}_i(t)\}_{i=1}^3$ satisfying, for each $i$,
    \begin{align*}
        m_i \ddot{\mathbf{r}}_i(t)
        =
        \sum_{\substack{j=1\\ j\neq i}}^{3}
        \frac{G m_i m_j\big(\mathbf{r}_j(t)-\mathbf{r}_i(t)\big)}
        {\|\mathbf{r}_j(t)-\mathbf{r}_i(t)\|^3},
        \qquad
        \text{whenever } \mathbf{r}_i(t)\neq \mathbf{r}_j(t).
    \end{align*}
    The set
    \[
        \Delta \defequal \big\{(\mathbf{r}_1,\mathbf{r}_2,\mathbf{r}_3)\in(\mathbb{R}^3)^3
        \mid \mathbf{r}_i=\mathbf{r}_j \text{ for some }i\neq j\big\}
    \]
    is called the \textit{collision set}.
\end{definition}

\begin{definition}[Phase Space and Hamiltonian]
    Put $\mathbf{v}_i \defequal \dot{\mathbf{r}}_i$ and consider the phase space
    \[
        \mathcal{P}\defequal \big((\mathbb{R}^3)^3\setminus \Delta\big)\times(\mathbb{R}^3)^3.
    \]
    Define the \textit{kinetic energy} $T$ and \textit{potential energy} $U$ by
    \begin{align*}
        T(\mathbf{v}) &\defequal \frac12\sum_{i=1}^3 m_i\|\mathbf{v}_i\|^2,\\
        U(\mathbf{r}) &\defequal -\sum_{1\le i<j\le 3}\frac{Gm_im_j}{\|\mathbf{r}_i-\mathbf{r}_j\|}.
    \end{align*}
    The \textit{total energy (Hamiltonian)} is
    \[
        H(\mathbf{r},\mathbf{v}) \defequal T(\mathbf{v})+U(\mathbf{r}).
    \]
\end{definition}

\begin{lemma}[Local Existence and Uniqueness]
    For any initial condition $(\mathbf{r}(0),\mathbf{v}(0))\in\mathcal{P}$,
    there exists $\varepsilon>0$ and a unique solution
    $(\mathbf{r}(t),\mathbf{v}(t))$ defined on $(-\varepsilon,\varepsilon)$.
\end{lemma}

\begin{definition}[Types of Collisions]
    A time $t_0$ is called a \textit{binary collision} if $\mathbf{r}_i(t_0)=\mathbf{r}_j(t_0)$ for exactly one pair $(i,j)$.
    It is called a \textit{triple collision} if $\mathbf{r}_1(t_0)=\mathbf{r}_2(t_0)=\mathbf{r}_3(t_0)$.
\end{definition}

\end{document}